\section{黏滞双重性：作为系统稳定器的阻尼}

\subsection{物理锚点：耗散结构的稳定性}

先看一个反常的事实。如果智能的扩张真如最激进的预测那样，是一场指数级爆发的海啸，那么在算力跨过某个临界点之后，物质世界应该已经被吞没。实际发生的是：比特领域的能力在飞速膨胀，而文明的物理实体与社会结构保持着惊人的黏滞。这种平静不能归因于人类拥有某种特权，也不能归因于技术停滞，原因更简单：文明的演进从来不在光滑的真空里进行。在比特世界的指数增长与原子世界的线性更新之间，立着两道阻尼。

这里需要区分执行与确认。执行受到资源、过程时间和既有结构约束，确认则依赖规则、对象证据与责任安排。两者共同影响技术部署的速度，也都可能被工程改进改变。由此形成的缓冲有多长，需要对具体系统判断。先用一个阻尼振子说明：外部驱动相同，惯性、阻尼与回复作用不同，响应就会不同。

\begin{equation}
m\frac{d^{2}x}{dt^{2}}+\gamma\frac{dx}{dt}+kx=F(t),
\qquad
\zeta=\frac{\gamma}{2\sqrt{mk}},
\qquad
\tau\approx\frac{\gamma}{k}\quad(\zeta\gg1)
\end{equation}

$x$ 为现实系统偏离原状态的程度，$m$、$\gamma$、$k$ 都是等效参量，若把 $x$ 取作长度，则 $m$ 为等效惯性，对应已固化的质量与沉没资本，单位千克；$\gamma$ 为等效阻尼系数，对应本书所说的资源、时间和结构三类黏滞，单位千克每秒；$k$ 为等效回复刚度，代表系统自我维持的能力，单位千克每二次方秒；$F(t)$ 为外部驱动力，即技术冲击，单位牛顿；$\zeta$ 为阻尼比，无量纲；$\tau$ 为过阻尼条件下的恢复时间，单位秒。当 $\zeta\gg1$ 时系统不产生振荡，而是以时间常数 $\tau\approx\gamma/k$ 单调趋近新的平衡，阻尼越大，恢复越慢。冲击的高度由 $F$ 决定，过渡的长度由 $\gamma/k$ 决定，两者可以分开计算：突变被拉成过渡期，不是因为冲击变小，而是因为阻尼变大。它的适用边界必须说清：这是一个线性、常系数的二阶模型，只用来给出阻尼与响应时间的方向关系；现实中的三类黏滞不是常数，它们会随技术输入重新分布，因此 $\tau$ 只能读成量级估计，不能用来预测具体的年份或进度。

维持运转需要持续付费。身体靠代谢，服务器靠供电与冷却，交通网靠巡检和维修；建筑可以在无人照料时保持相当长时间，也会在材料老化中逐步失去性能。这些结构的维持机制和时间尺度各不相同。一条旧路常常比一处空地难施工，因为它在被修补时还要承担交通。维护，是系统继续存在于使用之中的代价。

\subhead{资源势垒：必须支付的物理账单}

第一笔支出发生在物质转化的路径上。炼铁需要还原氧化铁，水泥制造需要驱动反应并处理高温物料。燃料、设备、反应与热流都要入账。活化势垒、反应自由能差和工厂实际能耗，是三种不同的量；催化、工艺和能源路线可以改变后两者的实现方式。智能的作用，正是在这些条件里找到更好的路，而每条路都要真的走过。

要有足够资源推动，这是资源势垒（$R$）。资源势垒指现实状态转移必须支付的能量、材料、设备、资本、空间、供应链和维护成本。信息可以描述一座桥，但桥本身必须由钢材、水泥、机械、土地和劳动建成；算法可以优化电网调度，但电力仍需要发电、输送、储能和调峰；AI 可以生成机器人动作策略，但机器人必须依靠电池、电机、传感器、关节、材料和维修体系行动。AI 可以降低设计成本，却不能取消物理账单。所谓“云”并不轻盈，它只是把原子负担隐藏在用户界面之后。越是走向大规模智能化，越依赖芯片、电力、数据中心、冷却系统、制造能力和基础设施。

人的身体也在这张账单上。每天摄入2500大卡，折算的全天平均代谢输入约为120瓦，要同时维持温控、认知、运动与修复。机器人与人比较，必须选相同任务、负载、工作质量和占空比，量出各自的平均输入。灵活的人在杂乱现场可能很经济，专用机器在标准化工序里也可能远远胜出。差别需要测量，不能从“碳基”或“硅基”直接读出。

这个换算本身就是那笔账单的算法，

\begin{equation}
P_{\text{人}}=\frac{E_{\text{食物}}}{\Delta t}
=\frac{2500\times4.18\times10^{3}}{86400}\approx1.2\times10^{2}\ \text{W},
\qquad
\frac{P_{\text{机}}}{P_{\text{人}}}=\frac{(1\text{--}3)\times10^{3}}{1.2\times10^{2}}\approx8\text{--}25
\end{equation}

$P_{\text{人}}$是摄入食物能折算的平均代谢功率，单位为瓦，$E_{\text{食物}}$是一天的食物能，$\Delta t$为一天的时长。以2500大卡计，平均功率约为120瓦，机械输出还要按活动与效率另算。$P_{\text{机}}$应采用同一任务下的实测平均输入功率；式中1到3千瓦只是比较示例，相应比值约为10到25。能量价格、设备成本与工作质量另行计入经济账本。

这个剪刀差还有一个容易被忽略的性质：它只规定价格，不规定方向。它保护的是劳动的相对成本优势，不是某项技能本身。一旦某个环节的任务被充分结构化，把人的那部分动作拆成一组可枚举的步骤，能耗优势就不再是决定性的，因为此时被比较的不再是通用智能和专用智能，而是两台专用设备。资源势垒的厚度在不同场景里差别很大，这正是它需要逐案计算而不能整段引用的原因。

\subhead{时间刚性：不可压缩的过程}

要有足够时间展开，这是时间刚性（$T$）。时间刚性指现实系统中不能任意压缩的过程时间。AI 可以瞬间生成建筑图纸，却不能让混凝土瞬间完成养护；可以快速设计分子，却不能让人体瞬间呈现长期副作用；可以模拟交通流，却不能让道路、车辆、法规、保险和驾驶习惯同步改变。物理系统有传热、扩散、疲劳和磨损；生命系统有生长、代谢、免疫和长期反馈；工程系统有测试、调试、认证和可靠性积累；制度系统有协商、合法化和行为适应。时间不是单纯的低效，有些慢是冗余，有些慢则是安全、信任和责任的边界。压缩时间往往不是免费加速，而是把风险推向未来。

时间刚性落在具体过程上。混凝土强度发展受配合比和养护条件影响；作物生长受品种、光照、温度和生理过程制约；核电站与高铁建设还要经过土地、环境、供应链和施工协调。基础设施的使用寿命也在规定更新节奏：住房和桥梁按几十年的周期设计\cite{miatto2017how}，改造必须与既有使用相衔接。并行施工、预制和更好的试验设计可以缩短部分周期，材料成熟、损伤暴露和利益协调则各有必要条件。时间刚性需要逐环节辨认。

这种脱节有一个直接的后果。AI 也许能把设计环节从一年压缩到一天，很难把施工环节从三年压缩到三个月。于是在相当长的时间里，人的头脑生活在 AI 编织的未来里，身体还寄居在缓慢更新的旧外壳中。制度层面的情况类似：一项技术的合法性来自协商和试错，而协商的时长不由技术效率决定，反而常常由需要被说服的人数决定。把时间刚性误读为低效，会导向一个反复出现的判断错误：以为把工具换得更快，整个流程就会跟着变快。真正会变快的是方案的数量，不是落地的速度。

\subhead{拓扑耦合：嵌入既有结构}

要能在既有结构中传导，这是拓扑耦合（$G$）。拓扑耦合指系统内部要素之间的连接方式、相互依赖、反馈回路，以及过去的选择对当前可选范围造成的限制。现实系统不是孤立变量的集合，而是由设备、人员、制度、习惯、标准、利益和环境构成的复杂网络。一个更聪明的自动驾驶算法，不等于交通系统自动优化；一个更好的医疗模型，不等于医院流程、医生责任、患者信任、医保支付和法律制度同步调整；一个更优的能源设备，不等于电网、储能、土地、资本和政策结构立刻适配。很多技术失败，不是因为单点能力不够，而是因为它无法嵌入既有系统的结构。

Moravec悖论把这种困难照得很清楚。规则明确的棋类可以拥有巨大的搜索空间，却仍然把状态、规则和胜负条件交代清楚。机器人处理一堆衣物，要同时应对遮挡、形变、接触与摩擦。人轻易完成的动作，背后已有一副适应环境的身体和长期形成的感觉运动能力。困难来自环境与任务的结构，不能单凭“高熵”两个字代替分析。

这个“指数级”有确切的来源，它来自构型空间的维度，

\begin{equation}
N_{\text{构型}}=m^{\,n},
\qquad
\ln N_{\text{构型}}=n\ln m
\end{equation}

式中$N_{\text{构型}}$是对$n$个自由度各取$m$个离散状态后得到的构型数，$n$、$m$和$N_{\text{构型}}$均为计数。构型数随自由度指数增长，说明逐点穷举很快会失去可行性。实际控制可以利用几何结构、学习、反馈和身体的被动力学缩小搜索范围，计算代价取决于采用的方法。抓取的困难还来自接触位置、摩擦和物体形变的不确定性。状态空间大小提供了一项困难来源，Moravec悖论的解释还需要演化经验、感知与控制机制。

物理世界的混乱确实保护了人，但它保护的是可被容忍的模糊性，不是所有劳动。现场判断之所以昂贵，是因为它必须在信息不完整的条件下做决定并承担后果，而不是因为它需要动手。一旦某个高熵场景被充分结构化，比如仓库被改造成标准货架，保护就随之下滑。这也说明 R、T、G 三者不是三面独立的墙：把一个场景整理干净，降低的是拓扑耦合，同时也在把这项劳动从资源势垒的保护下移出去。

\subhead{三种介质：比特、原子、生命}

不同介质对 AI 的吸收能力并不相同。比特、原子和生命，可以视为 AI 进入现实世界时面对的三种典型介质。三者不是静态分类，而是智能进入现实时逐级加重的介质梯度：比特更容易生成和回滚，原子必须承受材料、能量和维护，生命则进一步牵涉生长、代偿、长期后果和不可轻易撤销的干预。要看清梯度是怎么加重的，需要逐个看它们的吸收机制。

比特域包括文本、图像、代码、数据、模型和虚拟空间。它具有低复制成本、低传播成本、低修改成本和低回滚成本，因此最先被生成式 AI 改造。写作、翻译、摘要、搜索、客服、编程和图像生成之所以迅速爆发，是因为它们处在低 RTG 环境中。AI 的搜索与生成能力几乎可以直接转化为可见产出。但比特域并非零黏滞。它的低黏滞主要发生在生成层面，而不是可信层面。AI 可以低成本生成文本、图像、视频和代码，但社会真正需要的是可靠代码、真实图像、可验证知识和可追责判断。生成不是问题之后，验证成为问题；表达不是问题之后，判断成为问题；信息不是问题之后，信任成为问题。

原子域包括制造、建筑、交通、能源、物流、机器人、芯片、城市和基础设施。它必须面对质量、能量、摩擦、惯性、材料强度、空间约束和维护成本。原子世界不能像文本一样复制、删除、重写和回滚。AI 可以优化设计、调度和控制，却不能取消施工、运输、装配、调试、损耗和责任。具身智能正是原子域黏滞的集中体现。普通生成式 AI 主要在表示空间中运行，错误可以重生成；机器人则必须把表示变成行动，把预测变成控制，把目标变成受力、接触、运动和反馈。文本里的错误可以重生成，浇筑进构件里的错误只能凿掉重来。因此，聊天机器人可以迅速进入屏幕，而通用机器人、自动驾驶和智能基础设施进入现实的速度远慢于舆论想象。

生命域包括身体、细胞、药物、遗传、生态和生物反馈。它比一般原子系统更复杂，因为生命不是静态材料，而是会响应、代偿、适应、反弹甚至演化的复杂系统。AI 可以加速靶点发现、蛋白质预测、分子设计和实验筛选，但不能同步压缩毒理、药代、临床试验、长期安全性和个体差异观察。DNA 可以被写成序列，但基因功能并不等于文本意义；它必须在细胞环境、发育阶段、调控网络和外部生态中实现。生命不是软件，不能像软件一样轻易回滚。发现可以加速，验证却具有更强的时间刚性和非线性风险。

因此，AI 时代的基本图景不是均匀加速，而是不均匀吸收。比特流动，原子抵抗，生命反馈。AI 在低黏滞领域迅速扩散，在高黏滞领域缓慢嵌入，在生命和社会系统中则遭遇更强的验证、责任和信任约束。所谓“AI 改变世界”，并不是智能以同一速度穿透所有介质，而是智能在不同介质中被不同程度地吸收、阻滞、放大、扭曲和重新分配。

\begin{figure}[ht]
\centering
% >>>>>>>>>> figs/fig-media-gradient.tex >>>>>>>>>>
% 图 1-1  三种介质的吸收机制与特征时间尺度
% 主体是底部的 log-tau 尺度条；上方的 3x3 表只留 2-6 字短签。
\begin{tikzpicture}[x=1cm, y=1cm, font=\small,
  cell/.style={anchor=center, align=center, inner sep=1pt, font=\scriptsize, text=ink},
  rowlab/.style={anchor=east, align=right, font=\scriptsize\heiti, text=slate},
  headlab/.style={anchor=center, font=\small\heiti},
  tick/.style={anchor=north, font=\scriptsize, text=ink},
  barname/.style={anchor=west, font=\footnotesize\heiti, text=white},
  barspan/.style={anchor=east, font=\scriptsize, text=white}
]
\useasboundingbox (-1.58,0.56) rectangle (11.52,-5.00);

\def\dec{0.90}      % 每十倍程 0.90 cm（10^-3 -> 10^9 共 12 个十倍程）
\def\rh{0.54}       % 数据行高
\def\hdr{0.44}      % 表头高

% ================= 3 x 3 表（列用三色，行用底色深浅）=================
% 表头
\fill[deepblue!34] (0.00,\hdr) rectangle (3.60,0);
\fill[terra!34]    (3.60,\hdr) rectangle (7.20,0);
\fill[forest!34]   (7.20,\hdr) rectangle (10.80,0);
\node[headlab, text=deepblue!95] at (1.80,\hdr*0.5) {比特};
\node[headlab, text=terra!95]    at (5.40,\hdr*0.5) {原子};
\node[headlab, text=forest!95]   at (9.00,\hdr*0.5) {生命};
\draw[slate!55, line width=0.6pt] (0,\hdr) -- (10.80,\hdr);

% 行底色：同一列内越往下越浅
\foreach \i/\dep in {1/26, 2/15, 3/7} {
  \pgfmathsetmacro\yy{-\rh*(\i-1)}
  \pgfmathsetmacro\yb{-\rh*\i}
  \fill[deepblue!\dep] (0.00,\yy) rectangle (3.60,\yb);
  \fill[terra!\dep]    (3.60,\yy) rectangle (7.20,\yb);
  \fill[forest!\dep]   (7.20,\yy) rectangle (10.80,\yb);
  \draw[slate!20, line width=0.25pt] (0,\yb) -- (10.80,\yb);
}
\foreach \x in {3.60, 7.20} {
  \draw[slate!20, line width=0.25pt] (\x,\hdr) -- (\x,-1.62);
}
\draw[slate!55, line width=0.6pt] (0,-1.62) -- (10.80,-1.62);

% 行首短签
\node[rowlab] at (-0.14,-0.27) {吸收机制};
\node[rowlab] at (-0.14,-0.81) {主导约束};
\node[rowlab] at (-0.14,-1.35) {可逆性};

% 单元格：每格 2-6 字
\foreach \x/\a/\b in {%
  1.80/{扩散}/{注意力},%
  5.40/{做功}/{能量密度},%
  9.00/{反馈}/{代偿延迟}} {
  \node[cell, font=\small\heiti] at (\x,-0.27) {\a};
  \node[cell] at (\x,-0.81) {\b};
}
\node[cell, font=\small\heiti, text=deepblue!95] at (1.80,-1.35) {可回滚};
\node[cell, font=\small\heiti, text=terra!95]    at (5.40,-1.35) {有实损};
\node[cell, font=\small\heiti, text=forest!95]   at (9.00,-1.35) {不可逆};

% ================= 底部：特征时间尺度（对数轴，图的主体）=================
\def\yax{-2.35}
\draw[ink, line width=0.7pt] (-0.02,\yax) -- (10.82,\yax);

% 主刻度（每 3 个十倍程）与中间次刻度
\foreach \e in {-3,0,3,6,9} {
  \pgfmathsetmacro\xx{(\e+3)*\dec}
  \draw[ink, line width=0.7pt] (\xx,\yax) -- (\xx,\yax-0.11);
  \node[tick] at (\xx,\yax-0.13) {$10^{\e}$};
}
\foreach \e in {-2,-1,1,2,4,5,7,8} {
  \pgfmathsetmacro\xx{(\e+3)*\dec}
  \draw[slate!60, line width=0.4pt] (\xx,\yax) -- (\xx,\yax-0.06);
}
\node[anchor=north west, font=\scriptsize, text=slate] at (10.95,\yax-0.13) {秒};

% 三条载体带（起止与轴对齐）
% 比特：10^-3 - 10^1   -> x 0.00 - 3.60
\fill[deepblue!88] (0.00,-3.42) rectangle (3.60,-2.92);
\node[barname] at (0.14,-3.17) {比特};
\node[barspan] at (3.46,-3.17) {毫秒--十秒};
% 原子：10^0 - 10^7    -> x 2.70 - 9.00
\fill[terra!88] (2.70,-4.14) rectangle (9.00,-3.64);
\node[barname] at (2.84,-3.89) {原子};
\node[barspan] at (8.86,-3.89) {秒--月};
% 生命：10^4 - 10^9    -> x 6.30 - 10.80
\fill[forest!88] (6.30,-4.86) rectangle (10.80,-4.36);
\node[barname] at (6.44,-4.61) {生命};
\node[barspan] at (10.66,-4.61) {代际};

% 带端点引线：把带的起止投到轴上
\foreach \xx in {0.00,3.60,2.70,9.00,6.30,10.80} {
  \draw[slate!30, line width=0.3pt, dashed] (\xx,\yax) -- (\xx,\yax-0.11);
}
\end{tikzpicture}

% <<<<<<<<<< figs/fig-media-gradient.tex <<<<<<<<<<

\caption{三种介质的吸收机制与黏滞来源。比特域的低黏滞集中在生成层，可信层仍然昂贵；原子域每个环节都要支付物质与能量；生命域则把后果展开到时间与代际中。}
\label{fig:media-gradient}
\end{figure}

三者的高低不是固定的，介质的性质会随条件变化。文本一旦被用于决策就会带上责任，代码一旦进入生产系统就会带上回滚成本，药物候选一旦进入临床就会带上时间刚性。介质决定的是初始黏滞的等级，不是最终黏滞的数值。这也是为什么同一个模型能力，在不同场景里产生的震荡幅度差别极大。

R、T、G 三者并非孤立，而是相互置换。为了压缩时间，往往需要投入更多资源；为了节省资源，往往需要拉长时间；为了降低局部摩擦，往往需要增加接口、标准、冗余、监控、审计和合规结构。这个置换关系值得停一下。它意味着技术进步并不总是消灭阻力，更多时候是在资源、时间和结构之间重新分配阻力。一个系统在用户端变得轻盈，可能意味着后台变得更重；一个流程在表面上自动化，可能意味着整体结构更难维护。所以，\key{真正的问题不是某项技术是否“加速”，而是它降低了哪一类黏滞，又把黏滞转移到了哪里。}用高能耗换时间，用长时间换安全，用复杂度换局部效率，用冗余换可靠性，每一次置换都有代价，代价只是换了一个地方落地。

资源、时间和结构共同规定了现实改变的条件。技术可以降低总成本，也可以把某一环节的负担转到别处，两种结果需要分别核算。物理约束给出过程的下限，工程和制度决定实际代价如何发生、如何分担。执行条件得到满足之后，还有一个问题：生成出来的判断，凭什么可以被依赖。

\subsection{语义锚点：不完备性与系统的开放需求}

执行条件回答的是方案怎样进入现实，确认条件回答的是方案凭什么被依赖。格式完整、推演流畅和对象真实，分别要求不同的检验。模型可以参与这些检验，检验所依据的规则、资料、实验和责任安排也必须进入考察。验证的成本，正是在这些位置上分开。

\subhead{验证金字塔}

验证之所以成为瓶颈，是因为生成发生在表示空间，验证必须返回现实空间。现实空间有材料、能量、时间、噪声、复杂反馈和不可回滚后果。AI 可以生成一千个药物候选，但实验、临床和长期观察不能同等扩张；AI 可以提出无数驾驶策略，但真实道路中的天气、行人、施工、责任和长尾风险不能被一次模拟穷尽；AI 可以生成大量政策方案，但制度效果必须在组织、行为和公共信任中展开。验证不是单纯的认知审核，而是 RTG 的集中表现：验证需要资源，需要时间，需要覆盖复杂结构。

验证也不是一个单点动作，而是从形式正确到长期可承担的连续爬坡。可以把验证理解为一个金字塔。底层是形式验证，判断格式、语法、规则和局部逻辑是否成立，主要发生在比特域，黏滞较低。第二层是语义验证，判断含义、事实、引用和知识关系是否正确，开始涉及知识系统和来源追踪。第三层是因果验证，判断方案是否真的导致预期结果，而不只是看起来相关。第四层是工程验证，判断系统能否在真实噪声、故障、维护、成本和规模化条件下稳定运行。第五层是生命长期验证，判断药物、基因治疗、生态、教育或心理干预的长期后果是否可接受。最高层是社会责任验证，判断系统能否被解释、审计、追责、补偿并获得公共信任。把金字塔的形状写成数值，

\begin{equation}
N_{k}=N_{0}\prod_{i=1}^{k}p_{i},
\qquad
0\le p_{i}\le 1
\end{equation}

这里$N_0$为初始候选数，$p_i$为已通过前$i-1$层的候选在第$i$层的通过比例，$N_k$为通过前$k$层的候选数。连乘来自条件计数，不要求独立；若某层筛空，后续候选数保持为零。数量之外，每层还要支付检验成本。上层能接住多少候选，由这些成本与实际能力共同决定。

\begin{figure}[ht]
\centering
% >>>>>>>>>> figs/fig-verification-pyramid.tex >>>>>>>>>>
% 图 1-2  验证金字塔：六层，右侧用对数时间条代替原表的"典型周期"
% 原表内容压缩为：层宽=样本量、右侧条长=反馈周期、条右 4 字短签=主要失败模式。
\begin{tikzpicture}[x=1cm, y=1cm, font=\small,
  lay/.style={anchor=center, align=center, font=\scriptsize\heiti, text=ink, inner sep=0pt},
  tick/.style={anchor=north, font=\scriptsize, text=slate},
  fail/.style={anchor=west, font=\scriptsize\kaishu, text=rust!90}
]
% 左边原来切掉约 1.2pt：「代价递增」贴左缘。放 0.18cm。（用 _figclip.py 量出）
% 下缘原来收到 -1.16，(a)/(b) 小标题（y=-0.80 起、anchor=north）距页边只剩 0.56pt。
% 换成 Times 之后字形高度变了，余量被吃掉。放宽到 -1.45，留出约 12pt。
\useasboundingbox (-1.16,4.58) rectangle (11.96,-1.45);

\def\h{0.74}        % 层高
\def\hw{0.337}      % 每层半宽递变
\def\bx{5.86}       % 时间条起点
\def\decw{0.48}     % 每十倍程

% ================= 六层金字塔（半宽 2.70 -> 0.68）=================
% 颜色：冷 -> 暖，表示代价 / 周期上升
\foreach \k/\col/\dep in {1/deepblue/14, 2/deepblue/24, 3/slate/26,
                          4/terra/30, 5/rust/24, 6/rust/38} {
  \pgfmathsetmacro\ybot{\h*(\k-1)}
  \pgfmathsetmacro\ytop{\h*\k}
  \pgfmathsetmacro\wb{2.70-\hw*(\k-1)}
  \pgfmathsetmacro\wt{2.70-\hw*\k}
  \fill[\col!\dep] (2.70-\wb,\ybot) -- (2.70+\wb,\ybot) -- (2.70+\wt,\ytop) -- (2.70-\wt,\ytop) -- cycle;
  \draw[\col!55, line width=0.35pt]
        (2.70-\wb,\ybot) -- (2.70+\wb,\ybot) -- (2.70+\wt,\ytop) -- (2.70-\wt,\ytop) -- cycle;
}
\node[lay] at (2.70,0.37) {形式验证};
\node[lay] at (2.70,1.11) {语义验证};
\node[lay] at (2.70,1.85) {因果验证};
\node[lay] at (2.70,2.59) {工程验证};
\node[lay] at (2.70,3.33) {生命长期};
\node[lay] at (2.70,4.07) {社会责任};

% 左侧：代价递增
\draw[-{Stealth[length=2.2mm]}, slate!70, line width=0.6pt] (-0.48,-0.02) -- (-0.48,4.46);
\node[rotate=90, font=\scriptsize\heiti, text=slate, anchor=south] at (-0.74,2.22) {代价递增};

% ================= 右侧：反馈周期（对数条，同一基准）=================
% 轴：10^0 -> 10^10 秒
\def\yax{-0.22}
\draw[ink, line width=0.7pt] (\bx-0.04,\yax) -- (\bx+\decw*10+0.04,\yax);
\foreach \e in {0,2,4,6,8,10} {
  \pgfmathsetmacro\xx{\bx+\decw*\e}
  \draw[ink, line width=0.6pt] (\xx,\yax) -- (\xx,\yax-0.10);
  \node[tick] at (\xx,\yax-0.12) {$10^{\e}$};
}
\foreach \e in {1,3,5,7,9} {
  \pgfmathsetmacro\xx{\bx+\decw*\e}
  \draw[slate!55, line width=0.35pt] (\xx,\yax) -- (\xx,\yax-0.05);
  \draw[slate!18, line width=0.3pt, dashed] (\xx,\yax+0.04) -- (\xx,4.44);
}
\node[anchor=north west, font=\scriptsize, text=slate] at (\bx+\decw*10+0.32,\yax-0.12) {秒};

% 六根条 + 条右短签 + 连到金字塔的虚线
\foreach \k/\e/\col/\dep/\tag in {%
  1/2/{deepblue}/14/{假通过},%
  2/4/{deepblue}/24/{含义漂移},%
  3/6/{slate}/30/{误当因果},%
  4/8/{terra}/34/{现场失效},%
  5/9/{rust}/28/{方向反转},%
  6/10/{rust}/42/{代价外溢}} {
  \pgfmathsetmacro\yy{4.44-\h*(\k-0.5)}
  \pgfmathsetmacro\xe{\bx+\decw*\e}
  \pgfmathsetmacro\wr{2.70-\hw*(\k-0.5)}
  \fill[\col!\dep] (\bx,\yy+0.15) rectangle (\xe,\yy-0.15);
  \draw[\col!65, line width=0.4pt] (\bx,\yy+0.15) rectangle (\xe,\yy-0.15);
  \node[fail] at (\xe+0.16,\yy) {\tag};
  \draw[slate!28, line width=0.3pt, dashed] (2.70+\wr+0.03,\yy) -- (\bx-0.06,\yy);
}

% 层宽的含义
\node[anchor=north, font=\scriptsize\kaishu, text=slate]
      at (2.70,-0.16) {层宽＝有效样本量};

% ================= 栏标题（放在内容下方）=================
\node[anchor=north, font=\footnotesize\heiti, text=slate] at (2.70,-0.80) {(a) 验证金字塔};
\node[anchor=north, font=\footnotesize\heiti, text=slate] at (8.28,-0.80) {(b) 反馈周期};
\end{tikzpicture}

% <<<<<<<<<< figs/fig-verification-pyramid.tex <<<<<<<<<<

\caption{验证金字塔：从形式、语义、因果、工程、生命长期到社会责任，越往上越难自动化。}
\label{fig:verification-pyramid}
\end{figure}

形式检查、来源核对、实验和长期观察，可以以不同速度改进。自动检查工具降低了一部分验证成本，临床暴露、现场运行和责任处理仍要占用时间与资源。生成能力增长快于这些环节时，待验证候选会积压。瓶颈的位置由各层吞吐决定，验证金字塔把需要分别投资的环节列了出来。

金字塔的形状还会被另一种力量改变。越靠上的层，判断标准越依赖具体情境和承担主体，越难被写成一条可执行规则。形式验证可以定义成一套语法检查，语义验证可以定义成一套来源对照，社会责任验证却无法定义成任何一套静态规则，因为它的内容随制度、受害方和后果分布变化。把上层验证改写成下层可计算的问题，是工程上的常见做法，也是常见的失败方式：一旦某个指标被用来代替目标，被测对象就会围绕指标重新组织行为。测量本身进入了系统，这是后面要单独处理的问题。

\subhead{形式证明与确认依据}

形式系统能够证明什么，取决于它的公理与推理规则。哥德尔不完备性定理指出，一致、可有效公理化且足以表达基本算术的系统，存在自身无法判定的命题；在通常的可证明性条件下，它也不能在内部证明自身的一致性。\cite{godel1931} 定理描述了形式证明的边界。神经网络的学习与经验判断，还涉及数据、训练和对象反馈，需要按相应机制考察。

工程验证提供了一个更直接的区分：证明程序满足规范，仍要核对规范是否表达了实际要求。形式检查可以由机器完成，规范、传感器与测试场景也可以继续接受其他工具的检查；整条确认链最终依赖被选定的证据与标准。判断这条链是否足够，需要查明它覆盖了哪些错误，遗漏会造成什么后果。

人的位置由此进入制度层面。涉及利益和风险的标准，需要承受后果的人获得表达、授权和申诉的通道。人可以借助机器形成判断，判断的正当性还要由证据、程序和责任支持。形式能力与社会授权分别回答两个问题，这一区分将贯穿后面的主体与承责讨论。

\subhead{信任的锚点：可被惩罚的担保}

信任还依赖担保。Nassim Taleb以“切肤之痛”强调，决策者与后果承担者之间的利益关系会改变判断的可信程度。自然人可以用财产、职业和信誉承诺，组织也可以通过资产、身份和责任安排提供担保。需要检查的是：决定错误以后，损失能否回到作出决定的边界，纠正与赔偿能否实际发生。

医疗、法律和其他高后果决定，需要同时建立可靠性与追责机制。专业复核、记录、暂停权和责任安排，使决定可以被检查，也使错误有人处理。人工参与要配有真实权限和必要信息，自动执行则需要可追索的部署与担保结构。准确性与可追责性各有作用，不能以其中一项替代另一项。

信任还来自共同经历与持续关系。咨询、照护和陪伴中的理解，包含记忆、处境判断与承诺兑现。一次回应可以由模型生成，关系能否持续、危机由谁处理、错误如何纠正，则取决于服务背后的具体安排。真人参与的价值，需要在这些可辨认的关系和行动中说明。

确认与信任在这里汇合。一个判断需要可以追溯的依据，一项承诺需要能够被追责的承担者。验证可以自动化，担保可以制度化；两者都要有稳定的身份、记录和后果回流路径。人的身体、经历和有限性提供了具体的承受条件，制度则决定这些条件怎样进入判断与责任。

两道阻尼合起来，把原本可能瞬间发生的断裂拉伸成了一段过渡期。过渡期不是安稳地带。当两套时间表同时运行，最先出问题的不是设备，而是人对自身位置的判断：劳动市场按一种速度重排，自我评价按另一种速度塌陷。这种分裂不是心理问题，它有明确的结构，可以从替代发生的位置开始追。
